\section{LFU Cache}
\label{sec:sq-lfu-cache}

\problemstatement{Design a cache with a fixed capacity supporting
$\textit{get}(\textit{key})$ and $\textit{put}(\textit{key},
\textit{value})$, both in $O(1)$, where eviction (when the cache is full
and a new key must be inserted) removes the \textbf{least frequently
used} key; if multiple keys share the lowest frequency, the
\textbf{least recently used} among them is evicted.}

\begin{patternbox}
LFU is LRU (Section~\ref{sec:sq-lru-cache}) with a second tie-breaking
dimension added: eviction priority is now \emph{(frequency, recency)},
not recency alone. The keyword to notice: whenever a tie-breaking rule is
layered on top of a structure you already have working (here, LRU's
hash-map-plus-doubly-linked-list), try \textbf{bucketing} by the primary
key (frequency) and reusing the existing structure (an LRU list) as the
secondary tie-breaker \emph{within} each bucket, rather than designing an
entirely new structure from scratch.
\end{patternbox}

\subsection*{Understanding the problem carefully}

Group keys by their current access frequency into separate buckets. Each
bucket, internally, only ever needs to break ties by recency among keys
that share that exact frequency -- which is exactly the LRU Cache
substructure from Section~\ref{sec:sq-lru-cache}, reused once per
distinct frequency value. Eviction always targets the \emph{lowest}
currently-occupied frequency bucket, and within it, removes that
bucket's least-recently-used key (the LRU substructure's own eviction
rule, applied locally).

\begin{ideabox}[A Hash Map of Frequency Buckets, Each an LRU List]
Maintain: a hash map $\textit{keyData}$ from key to $(\textit{value},
\textit{freq})$; a hash map $\textit{freqBuckets}$ from frequency to a
doubly linked list (an LRU list) of keys currently at that frequency,
ordered by recency within the bucket; and a single integer
$\textit{minFreq}$, tracking the smallest frequency currently occupied by
any key. On every access to an existing key (via \textit{get} or an
updating \textit{put}): remove it from its current frequency bucket
(inserting it at the front, i.e.\ most-recent position, of the
\emph{next} frequency bucket up), incrementing its stored frequency; if
its old bucket becomes empty and was the minimum, increment
$\textit{minFreq}$. On inserting a brand-new key when the cache is full:
evict the least-recently-used key from the $\textit{minFreq}$ bucket, then
insert the new key into the frequency-$1$ bucket and reset
$\textit{minFreq} \gets 1$.
\end{ideabox}

\subsection*{Why tracking \textit{minFreq} explicitly avoids an expensive search}

Without an explicit running $\textit{minFreq}$, finding ``the lowest
currently-occupied frequency'' would require scanning all buckets,
costing more than $O(1)$. Because frequency only ever \textbf{increases}
by exactly $1$ per access, and a brand-new key always enters at frequency
$1$, we can maintain $\textit{minFreq}$ incrementally with simple,
local updates: it can only decrease (to $1$) when a new key is inserted,
and it can only increase when the bucket it currently points to becomes
completely empty (because the key that just left was the last one at that
frequency, and every other key's frequency is, by definition, at least
\textit{minFreq}, so the new minimum must be \textit{minFreq}+1 in that
specific case). This local bookkeeping replaces a potentially expensive
global search with $O(1)$ updates.

\subsection*{Algorithm}

\begin{enumerate}
  \item \textsc{Touch}($\textit{key}$) (shared helper for get and
        updating put): look up $(\textit{value}, f)$ in
        $\textit{keyData}$; remove $\textit{key}$ from
        $\textit{freqBuckets}[f]$; if that bucket is now empty and
        $f = \textit{minFreq}$, increment \textit{minFreq}; insert
        $\textit{key}$ at the front of $\textit{freqBuckets}[f+1]$;
        update $\textit{keyData}[\textit{key}] \gets (\textit{value},
        f+1)$.
  \item \textsc{Get}($\textit{key}$): if absent, return $-1$. Else:
        \textsc{Touch}($\textit{key}$); return its value.
  \item \textsc{Put}($\textit{key}, \textit{value}$): if present, update
        the stored value and call \textsc{Touch}($\textit{key}$).
        Otherwise, if at capacity: evict the least-recently-used key
        from $\textit{freqBuckets}[\textit{minFreq}]$ (and its
        $\textit{keyData}$ entry). Insert $\textit{key}$ at the front of
        $\textit{freqBuckets}[1]$; set $\textit{keyData}[\textit{key}]
        \gets (\textit{value}, 1)$; set $\textit{minFreq} \gets 1$.
\end{enumerate}

\subsection*{Dry run}

\dryrun{Capacity $2$. \textsc{Put}(1,1), \textsc{Put}(2,2),
\textsc{Get}(1) (freq of $1$ becomes $2$), \textsc{Put}(3,3) (cache full;
evict lowest-frequency key, which is $2$, frequency $1$), \textsc{Get}(2).}

\begin{itemize}
  \item Put(1,1): bucket[1]$=[1]$. minFreq$=1$.
  \item Put(2,2): bucket[1]$=[2,1]$ (2 most recent). minFreq$=1$.
  \item Get(1): touch $1$ (freq $1\to2$). Remove from bucket[1]:
        bucket[1]$=[2]$ (not empty, minFreq stays $1$). Insert into
        bucket[2]$=[1]$.
  \item Put(3,3): cache full ($2$ keys). Evict from bucket[\textit{minFreq}=1]'s
        LRU end: bucket[1]$=[2]$, evict $2$. bucket[1] now empty. Insert
        $3$ into bucket[1]$=[3]$. minFreq reset to $1$.
  \item Get(2): key $2$ was evicted; returns $-1$.
\end{itemize}

\subsection*{C++ implementation}

\begin{lstlisting}
#include <bits/stdc++.h>
using namespace std;

class LFUCache {
    int capacity, minFreq;
    unordered_map<int, pair<int,int>> keyData; // key -> (value, freq)
    unordered_map<int, list<int>> freqBuckets;  // freq -> keys (front = most recent)
    unordered_map<int, list<int>::iterator> keyIter; // key -> its position in its bucket

    void touch(int key) {
        auto [value, freq] = keyData[key];
        freqBuckets[freq].erase(keyIter[key]);
        if (freqBuckets[freq].empty()) {
            freqBuckets.erase(freq);
            if (freq == minFreq) minFreq++;
        }
        freqBuckets[freq + 1].push_front(key);
        keyIter[key] = freqBuckets[freq + 1].begin();
        keyData[key] = {value, freq + 1};
    }

public:
    LFUCache(int cap) : capacity(cap), minFreq(0) {}

    int get(int key) {
        if (!keyData.count(key)) return -1;
        touch(key);
        return keyData[key].first;
    }

    void put(int key, int value) {
        if (capacity == 0) return;

        if (keyData.count(key)) {
            keyData[key].first = value;
            touch(key);
            return;
        }

        if ((int)keyData.size() == capacity) {
            int evictKey = freqBuckets[minFreq].back();
            freqBuckets[minFreq].pop_back();
            keyData.erase(evictKey);
            keyIter.erase(evictKey);
        }

        keyData[key] = {value, 1};
        freqBuckets[1].push_front(key);
        keyIter[key] = freqBuckets[1].begin();
        minFreq = 1;
    }
};

int main() {
    LFUCache cache(2);
    cache.put(1, 1);
    cache.put(2, 2);
    cout << cache.get(1) << endl; // 1
    cache.put(3, 3);              // evicts key 2
    cout << cache.get(2) << endl; // -1
    return 0;
}
\end{lstlisting}

\begin{complexitybox}
\textbf{Time Complexity.} $O(1)$ for both $\textit{get}$ and
$\textit{put}$ -- every step is a hash map lookup or a constant number of
doubly linked list operations, with $\textit{minFreq}$ maintained
incrementally rather than recomputed.
\textbf{Space Complexity.} $O(\textit{capacity})$ across the hash maps
and linked lists.
\end{complexitybox}

\begin{takeawaybox}
This closing problem of the chapter is a deliberate capstone, much like
Median of a Row-wise Sorted Matrix was for the Binary Search chapter: it
requires recognizing that a new, seemingly more complex requirement
(frequency-based eviction) can be built by \emph{bucketing} on the new
dimension and reusing an already-solved structure (LRU's doubly linked
list) as the tie-breaker within each bucket -- composition of known
pieces, rather than a wholly new design, solves the harder problem.
\end{takeawaybox}
